\section{Toy Generation and Validation Studies}
\label{sec:toys}

\subsection{Scope of the current validation campaign}

Unless stated otherwise, the toy studies summarized here correspond to the 95\%
confidence-level workflow with blind width $1.64\,\sigma_m$ and \num{10000} toys per
mass-strength point. The completed suites cover the fully unblinded 2015 sample, the
2016 10\% development sample, and the first 2015+2016 shared-$\eps^2$ combination.
Throughout this section, the combined interpretation is focused on the overlap region
\SIrange{35}{130}{MeV}, which is the mass interval where both data sets contribute
meaningfully to the prompt-resonance search.

The bundled toy-validation figures were produced before the final adoption of the
leakage-retaining signal template and before the future likelihood refinements
mentioned in Section~\ref{sec:methodology}. They are nevertheless already informative
for the statistical questions that matter most: closure, bias, pull calibration,
coverage, and the logic of the combined extraction.

\subsection{Three toy-generation modes}

The validation program uses three complementary pseudoexperiment constructions.

\subsubsection{Analytic functional-form pseudoexperiments}

The first class of toys starts from an analytic fit to the smooth $e^+e^-$
invariant-mass spectrum and fluctuates pseudo-data around that closed-form background
model. These studies isolate the question of gross spectral smoothness from the later
question of GP conditioning. They are therefore useful as the first closure test before
the full inference chain is exercised.

\begin{figure}[H]
\centering
\begin{subfigure}[t]{0.32\linewidth}
  \centering
  \graphicorplaceholder{\linewidth}{toy_generation_figs/good_fit_2021_0pt03_pct.png}
  \caption{2021 0.03\% analytic fit.}
\end{subfigure}
\hfill
\begin{subfigure}[t]{0.32\linewidth}
  \centering
  \graphicorplaceholder{\linewidth}{toy_generation_figs/good_fit_2015_placeholder.png}
  \caption{2015 placeholder.}
\end{subfigure}
\hfill
\begin{subfigure}[t]{0.32\linewidth}
  \centering
  \graphicorplaceholder{\linewidth}{toy_generation_figs/good_fit_2016_placeholder.png}
  \caption{2016 placeholder.}
\end{subfigure}
\caption{Analytic background fits used to seed functional-form pseudoexperiments. The
2021 panel is already available; analogous 2015 and 2016 plots should be added when
those studies are exported in final form.}
\label{fig:analytic-fit-toys}
\end{figure}

\subsubsection{Conditional blind-window GP toys}

The fastest toy mode keeps the sideband-trained GP fixed at a given test mass. The GP
fit is performed once on the sidebands, the blind-window mean vector and covariance are
constructed from that fit, a background realization is drawn, signal is added, and the
result is Poisson-sampled. This construction is ideal for high-statistics expected-band
studies and for testing the response of the extraction fit under the posterior
background model preferred by the real data.

\subsubsection{Full procedural GP-mean/refit pseudoexperiments}

The most realistic toy mode generates a full pseudoexperiment from the GP-derived
background expectation, injects signal with the full Gaussian template, and then reruns
the GP sideband-to-blind prediction on the toy before signal extraction. This is the
appropriate mode for diagnosing signal absorption, because it directly tests how much
of the injected signal may be absorbed by the refitted background model or by the
profiled nuisance treatment.

\subsection{Asimov versus Poisson reference construction}

When toy strengths are specified in units of $\sigma_A$, the package first constructs a
reference extraction uncertainty $\sigma_{A,\mathrm{ref}}(m)$. In \textbf{Asimov mode},
the reference spectrum is set equal to the predicted background mean and therefore
encodes the deterministic sensitivity implied by the fit. In \textbf{Poisson mode}, the
reference spectrum is fluctuated before $\sigma_{A,\mathrm{ref}}$ is extracted. The
completed 95\% CL studies summarized here use the Asimov convention.

\subsection{Poisson versus multinomial signal injection}

Signal injection is distinct from the choice of $\sigma_A$ reference. In
\textbf{Poisson-total injection}, each bin receives a Poisson draw with mean
$A_{\mathrm{inj}}w_i$, so the realized total signal fluctuates from toy to toy. In
\textbf{multinomial fixed-total injection}, the total number of injected events is
fixed and only the bin-to-bin allocation fluctuates. Since the present signal template
is normalized on the full histogram range, $A_{\mathrm{inj}}$ denotes the total local
signal yield and only the fraction $f_{\mathrm{blind}}A_{\mathrm{inj}}$ is expected to
appear inside the blinded bins for the nominal $1.64\,\sigma_m$ window.

This distinction matters when interpreting uncertainty bars on validation plots. Only
Poisson-total injection produces a genuine horizontal spread in the realized injected
strength. Under fixed-total multinomial injection the horizontal coordinate is exact.

\subsection{Representative single-mass extraction displays}

Representative extraction displays make the statistical procedure concrete at the level
of an individual mass hypothesis. The 2015 \SI{30}{MeV} display and the 2016 10\%
\SI{58}{MeV} display are retained because they show, in real data, how the profiled
background deformation competes with the narrow signal template. The 2016 example is
particularly useful because the best-fit amplitude is negative.

\begin{figure}[p]
\centering
\begin{subfigure}[t]{0.88\linewidth}
  \centering
  \graphicorplaceholder{\linewidth}{fit_example_figs/2015/m030MeV_blind_fit.png}
  \caption{2015 data at \SI{30}{MeV}.}
\end{subfigure}

\vspace{1.0em}

\begin{subfigure}[t]{0.88\linewidth}
  \centering
  \graphicorplaceholder{\linewidth}{fit_example_figs/2016/m058MeV_blind_fit.png}
  \caption{2016 10\% data at \SI{58}{MeV}, including a negative $\hat A$.}
\end{subfigure}
\caption{Representative blind-window fits using data. These displays show the fit
components directly and serve as concrete examples of the local extraction machinery
before one moves to ensemble diagnostics.}
\label{fig:data-fit-examples}
\end{figure}

\subsection{Pull, coverage, and significance-calibration diagnostics}

The central ensemble-level validation quantities are the pull,
\begin{equation}
\mathrm{pull} = \frac{\hat A - A_{\mathrm{inj}}}{\sigma_A},
\end{equation}
its mean and width as a function of mass and injected strength, and the empirical
coverage of the nominal one- and two-standard-deviation intervals. These are standard
closure diagnostics in HEP statistical analyses because they test, respectively,
whether the estimator is biased, whether the quoted uncertainty is correctly calibrated,
and whether the reported intervals achieve the claimed frequentist coverage
\cite{Cowan1998,Cowan2011}.

For the present study, the pull-versus-mass plots are the most compact summary because
they display both the pull mean and pull width at once. In the overlap region
\SIrange{40}{115}{MeV}, both the 2015 and 2016 10\% studies show broadly promising
behavior. The large downward excursion below \SI{40}{MeV} is dominated by the fact that
the 2016 geometric acceptance is essentially zero there, especially at
\SIlist{25;30}{MeV}. The upward drift above approximately \SI{130}{MeV} is influenced by
the truncation of the 2015 spectrum near \SI{150}{MeV} and may also contain residual
signal-leakage effects that should be revisited after the next likelihood update.

\begin{figure}[p]
\centering
\begin{subfigure}[t]{0.48\linewidth}
  \centering
  \graphicorplaceholder{\linewidth}{injection_extraction_figs/pull_vs_mass_2015.png}
  \caption{2015 pull mean and width versus mass.}
\end{subfigure}
\hfill
\begin{subfigure}[t]{0.48\linewidth}
  \centering
  \graphicorplaceholder{\linewidth}{injection_extraction_figs/pull_vs_mass_2016.png}
  \caption{2016 10\% pull mean and width versus mass.}
\end{subfigure}
\caption{Mass-dependent pull diagnostics for the single-dataset analyses. The points
show the mean pull and pull width extracted from the toy ensembles, while the shaded
bands reflect the uncertainty of those quantities under the Poisson injection scheme.
These plots summarize the most important closure information more directly than the
separate linearity, bias, and pull-width aggregate summaries.}
\label{fig:pull-vs-mass}
\end{figure}

The coverage plots complement the pull diagnostics by testing whether the empirical
fractions satisfying $|\mathrm{pull}|<1$ and $|\mathrm{pull}|<2$ remain close to the
nominal Gaussian expectations of approximately 68\% and 95\%. For both 2015 and the
2016 10\% sample, the observed coverage is reasonably close to expectation across the
core search region.

\begin{figure}[p]
\centering
\begin{subfigure}[t]{0.48\linewidth}
  \centering
  \graphicorplaceholder{\linewidth}{injection_extraction_figs/coverage_2015.png}
  \caption{2015 interval coverage.}
\end{subfigure}
\hfill
\begin{subfigure}[t]{0.48\linewidth}
  \centering
  \graphicorplaceholder{\linewidth}{injection_extraction_figs/coverage_2016.png}
  \caption{2016 10\% interval coverage.}
\end{subfigure}
\caption{Empirical coverage of the nominal one- and two-standard-deviation pull
intervals. These panels test whether the extracted uncertainty $\sigma_A$ is correctly
calibrated as an interval statement, not only as a width parameter.}
\label{fig:coverage}
\end{figure}

The significance-calibration studies provide a third layer of validation by comparing
the extracted significance scale with the injected scale. This information corroborates
the combined-search-power studies discussed below: both are asking whether the
multi-dataset procedure accumulates significance in a way consistent with the expected
inverse-variance weighting of the constituent channels.

\subsection{Combined search-power studies}

The combined-search-power figures illustrate why a shared-$\eps^2$ combination can
outperform either constituent data set alone. The key point is that the combined
injection is not defined by adding the same number of signal events to each data set.
Instead, a common injected $\eps^2$ is mapped into dataset-specific total signal yields
through the conversion factors $K_d(m)$ and then interpreted relative to the
dataset-specific $\sigma_{A,\mathrm{ref}}$. As a result, a combined ``7$\sigma$''
injection means a seven-standard-deviation common coupling in the combined statistic,
not seven standard deviations in each constituent data set separately.

\begin{figure}[H]
\centering
\begin{subfigure}[t]{0.48\linewidth}
  \centering
  \graphicorplaceholder{\linewidth}{combined_search_figs/combined_search_power_scenarios.png}
  \caption{Scenario-level combined significance.}
\end{subfigure}
\hfill
\begin{subfigure}[t]{0.48\linewidth}
  \centering
  \graphicorplaceholder{\linewidth}{combined_search_figs/combined_search_power_constituent_pvalues_5sigma.png}
  \caption{Constituent requirements for a target combined $5\sigma$ excess.}
\end{subfigure}
\caption{Combined-search-power diagnostics. These plots show how sensitivity accumulates
in the shared-$\eps^2$ framework and should be interpreted together with the
significance-calibration studies rather than as direct limit-setting observables.}
\label{fig:combined-search-power}
\end{figure}

\subsection{Combined extraction displays}

The common-signal extraction displays are the most transparent way to visualize the
shared-$\eps^2$ construction. The top row shows individual single-dataset
pseudoexperiments. The lower panel shows the combined Gaussian signal representation
constructed from the same injected coupling. For the currently available \SI{40}{MeV}
example, the injected combined strength is seven standard deviations in the
\emph{combined} statistic, which corresponds to different effective per-dataset
strengths because the two data sets have different conversion factors and reference
uncertainties.

\begin{figure}[p]
\centering
\graphicorplaceholder{0.96\linewidth}{combined_search_figs/extract_display_combined_m040MeV_z7p0.png}
\caption{Illustrative combined extraction display at \SI{40}{MeV}. The annotation box
summarizes the injected coupling, the extracted combined coupling, the observed combined
limit, and the corresponding per-dataset fitted yields. The upper panels show the
single-dataset fits and extracted signal components, while the lower panel shows the
combined Gaussian signal representation that enters the shared-$\eps^2$ interpretation.
The visible negative pull is why this figure should still be regarded as a diagnostic
display rather than as the final combined-search illustration.}
\label{fig:combined-extract-displays}
\end{figure}
